\section*{Termocuplas por tipos}
Las termocuplas cuenta con una clasificación basado en letras según el tipo de material empleado para su creación, siendo así además que dentro de sus características se toman en cuenta otras consideraciones para el voltaje, sensibilidad a @25°, siendo el voltaje generado relativo a una unión de referencia a 0°C. El valor indicado es el generado cuando la unión de medida está a 25°C.
Mientras que la sensibilidad \textbf{coeficiente Seebeck} es aproximada a esa temperatura y varía a lo largo del rango. Para el tipo B, la salida es casi cero a 25°C, por lo que su sensibilidad se suele dar a temperaturas más altas.

Y dentro de las aplicaciones para cada tipo de termocupla, por ejemplo, se tiene que los tipos B, R y S son nobles, estables y para altísimas temperaturas (hornos) y por el contrario los tipos K, J, E y T son de metales base, más económicos y de uso general. El tipo N es una mejora moderna del tipo K, con mayor estabilidad.

% Please add the following required packages to your document preamble:
% \usepackage{multirow}
\begin{table}[h]
	\begin{tabular}{|c|ll|c|c|c|}
		\hline
		\multirow{2}{*}{\textbf{\begin{tabular}[c]{@{}c@{}}Tipo\\ ANSI\end{tabular}}} & \multicolumn{2}{c|}{\textbf{Composición}}                                                                                                 & \multirow{2}{*}{\textbf{\begin{tabular}[c]{@{}c@{}}Rango de\\ T (°C)\end{tabular}}} & \multirow{2}{*}{\textbf{\begin{tabular}[c]{@{}c@{}}Voltaje (mV)\\ @25°C\end{tabular}}} & \multirow{2}{*}{\textbf{\begin{tabular}[c]{@{}c@{}}Sensibilidad (uV/°C)\\ @25°C\end{tabular}}} \\ \cline{2-3}
		& \multicolumn{1}{c|}{\textbf{Composición Positiva}}                                               & \multicolumn{1}{c|}{\textbf{Negativo}} &                                                                                     &                                                                                        &                                                                                                \\ \hline
		\textbf{B}                                                                    & \multicolumn{1}{l|}{\begin{tabular}[c]{@{}l@{}}Platino-Rodio\\ (70\% Pt, 30\% Rh)\end{tabular}}  & Platino-Rodio                          & 0 a +1820                                                                           & $\sim$0,000                                                                            & \begin{tabular}[c]{@{}c@{}}Aprox. 8\\ (a 1000°C)\end{tabular}                                  \\ \hline
		\textbf{E}                                                                    & \multicolumn{1}{l|}{Níquel-Cromo (Cromel)}                                                       & Níquel-Silicio                         & -270 a +1000                                                                        & 1.495                                                                                  & $\sim$61                                                                                       \\ \hline
		\textbf{J}                                                                    & \multicolumn{1}{l|}{Hierro (Fe)}                                                                 & Níquel-Silicio                         & -210 a +1200                                                                        & 1.277                                                                                  & $\sim$51                                                                                       \\ \hline
		\textbf{K}                                                                    & \multicolumn{1}{l|}{Níquel-Cromo (Cromel)}                                                       & Níquel-Aluminio                        & -270 a +1372                                                                        & 1.000                                                                                  & $\sim$41                                                                                       \\ \hline
		\textbf{N}                                                                    & \multicolumn{1}{l|}{\begin{tabular}[c]{@{}l@{}}Níquel-Cromo-Silicio \\ (Nicrosil)\end{tabular}}  & NI-SI-MG                               & -270 a +1300                                                                        & 0,261                                                                                  & $\sim$30                                                                                       \\ \hline
		\textbf{R}                                                                    & \multicolumn{1}{l|}{\begin{tabular}[c]{@{}l@{}}Platino-Rodio \\ (87\% Pt, 13\% Rh)\end{tabular}} & Platino (Pt)                           & -50 a +1768                                                                         & 0,226                                                                                  & $\sim$12                                                                                       \\ \hline
		\textbf{S}                                                                    & \multicolumn{1}{l|}{\begin{tabular}[c]{@{}l@{}}Platino-Rodio \\ (90\% Pt, 10\% Rh)\end{tabular}} & Platino (Pt)                           & -50 a +1768                                                                         & 0,142                                                                                  & $\sim$10                                                                                       \\ \hline
		\textbf{T}                                                                    & \multicolumn{1}{l|}{Cobre (Cu)}                                                                  & Níquel-Silicio                         & -270 a +400                                                                         & 0,992                                                                                  & $\sim$41                                                                                       \\ \hline
	\end{tabular}
\end{table}

\section*{Principio de funcionamiento y partes}

Principio de funcionamiento: El principio fundamental de una termocupla es el efecto Seebeck, un fenómeno termoeléctrico descubierto en 1821. Este efecto establece que cuando dos metales o semiconductores de materiales diferentes se unen en sus extremos formando un circuito, y se mantiene una diferencia de temperatura entre la unión de medida (o unión caliente) y la unión de referencia (o unión fría), se genera una fuerza electromotriz (FEM) o voltaje eléctrico muy pequeño (del orden de milivoltios) en el circuito. Este voltaje es directamente proporcional a la diferencia de temperatura entre ambas uniones. Para determinar la temperatura absoluta en el punto de medición, es esencial conocer o compensar con precisión la temperatura en la unión de referencia, lo que se logra mediante circuitos electrónicos de compensación de unión fría en el instrumento lector.

\subsection*{Partes}
Partes principales (estructura básica) de una termocupla consta de dos elementos clave:

\begin{itemize}
	\item \textbf{Los pares de alambre:} Son dos conductores de metales distintos (por ejemplo, Cromel y Alumel para el tipo K), que constituyen el par termoeléctrico. Su composición define el rango, sensibilidad y características del sensor.
	
	\item \textbf{Las uniones:} La unión de medida (hot junction) es el punto donde se sueldan o fusionan los dos alambres y se pone en contacto con el medio cuya temperatura se desea medir. La unión de referencia (cold junction) es el punto donde los alambres se conectan al instrumento de medición (o donde se conoce su temperatura).
\end{itemize}